\section{Discussion}\label{sec:discussion}

\paragraph{Which failure signals predict successful recovery?} Three, in descending order of
strength. \emph{Class} is the dominant predictor, and specifically whether the class has a cheap
local operator: misuse and contradictory evidence are repaired at roughly 96\%, while planning error
(41\%) and objective drift (52\%) are not, because their repair requires reconstructing intent that
is not present in any observation. \emph{Timing} is second: detection before propagation is worth
roughly \EarlyGap{} points of recovery within the same policy. \emph{Signal type} is third and is
what makes class prediction possible at all---error codes alone detect 64\% of harmful faults, and
adding staleness, conflict, policy-conformance and progress features recovers the rest. Notably, the
predictors that matter are \emph{structural}, not stylistic: token counts, latency and retry counts
carry almost no information about whether a repair will succeed.

\paragraph{Repair has a ceiling that detection cannot lift.} The most useful negative result here is
the shape of the failure: detection AUC near 0.8 with recovery at 61\%. Even a perfect detector would
not close the gap, because the remaining failures are class-confusion and operator-locality problems.
This reframes the research agenda. Work that improves monitors (better step-level classifiers,
process reward models, LLM critics) attacks the smaller half of the problem. The larger half is
\emph{repair operators}: re-anchoring an objective, or replanning from a damaged belief state, is a
different capability from noticing that something is wrong, and it is where the numbers are low.

The LLM-controller experiment points the same way. The classes it repairs as well as the rules are
exactly the ones whose operator is named after their symptom (misuse, staleness); the classes whose
repair requires reconstructing intent collapse to retries. Operator locality, not signal visibility,
is again what separates repair from detection.

\paragraph{The learned-policy result is a data problem, not a modelling problem.} Our ranker failed
to beat the canonical mapping. We could have hidden this by tuning the candidate set, and we chose
not to, because the diagnosis is the interesting part: recovery policy learning from logged
trajectories is an off-policy problem, and the logging policy never takes the actions we need to
evaluate. Fixing it requires exploration in the sandbox (a small amount of deliberate operator
randomisation during data collection) rather than a better model class. That is a concrete,
actionable design for the next version of the benchmark, and it is exactly the kind of finding a
benchmark should produce about itself.

\paragraph{Why state-verified scoring is not optional.} The self-report gap (\ClaimedPct\% claimed
vs \VerifiedPct\% verified) is the clearest practical lesson. It is not that agents lie; it is that
the information needed to know they failed is not in the transcript. In the migration family a lost
update is invisible to an agent that only sees what it wrote. Any evaluation based on transcript
inspection, LLM-judge grading, or self-reported completion is measuring confidence, not reliability.

\paragraph{What a leaderboard should and should not reward.} \arb{} is designed so that a system
cannot win by being conservative (containment is scored separately and counts as failure to
recover), by being verbose (tokens are costed), by retrying (retry-only is a baseline at 0\%), or by
claiming success (claims are ignored). What it rewards is conditional recovery at bounded cost with
containment on the unrecoverable class. Because the harmful set is fixed by the screen, every
entrant is scored on the same faults, which is what makes cross-system comparison meaningful.

\paragraph{Limitations and threats to validity}\label{sec:limits}

\begin{itemize}
\item \textbf{Simulation.} Every environment is a deterministic Python world. Real tools are
partially observable, non-deterministic, and slow; real failures are heavily correlated with load,
version skew and human behaviour. \arb{} measures \emph{relative} reliability under controlled
faults. It cannot tell you an agent's production recovery rate, and a system that ranks first here
may not rank first in deployment.
\item \textbf{Scripted agents, LLM controller.} The agents that execute tasks are scripted workflow
executors with a known SOP. They are a strong control (deterministic, belief-grounded, replayable) and
a narrow sample of agent behaviour. We supplement them with a local 7B model used as the
\emph{recovery controller} (\LLMRecovery\% versus \LLMRuleRecovery\% for the rules on the same
pairs), which is a capability result about operator selection, not about planning. A frontier model
driving both the plan and the repair is still not evaluated, so the numbers characterise the
\emph{harness and the fault set} rather than any deployed model.
\item \textbf{Harm screening induces selection.} Excluding inert faults biases the retained set
toward faults with large downstream reach, which inflates apparent harmfulness and may make
recovery look harder than it is for the average injection. We chose this because the alternative---a
denominator containing faults that no monitor could be blamed for---makes recall uninterpretable.
\item \textbf{Cluster bootstrap over templates.} With a few dozen templates the bootstrap is
coarse; intervals should be read as indicating order of magnitude, not third-decimal precision.
\item \textbf{Single lab, single author.} Labels are defined by an injection program, which removes
annotator disagreement but also removes the chance to discover classes we did not think of. Natural
faults partially mitigate this, and their class agreement (\NaturalAgreement\%) is encouraging, but
the taxonomy is not exhaustive.
\item \textbf{Cost model.} The price sheet is a fixed artefact; USD figures are a transparent
normalisation, not a market forecast.
\end{itemize}

\paragraph{Ethics and safety.} Fault injection is confined to sandboxed, non-production
environments with no external side effects; no real refunds, orders or migrations occur. The
benchmark's emphasis on containment is deliberate: an agent that stops mutating when it does not
understand the world is safer, and scoring containment separately prevents the benchmark from
rewarding agents that push through uncertainty. Releasing a recovery policy that is only evaluated
in sandbox is a real risk---it invites deployment outside the tested boundary---so the released
ranker is packaged with an explicit out-of-scope statement, and the leaderboard protocol requires
sandbox evaluation of any submitted policy. We do not release attack tooling: the injection
programs perturb only the benchmark's own environments.
